% GENERATED FILE. Edit canonical lessons, then run npm run generate:books.
% canonical-source-hash: fnv1a64:5f76fd9e0cddb2de
% canonical-lessons: HI-C406-taiksi, HI-C406-tayar, HI-C406-pancar, HI-C406-helmet, HI-C406-laisens

\chapter{Taxis and Tyres}
\label{ch:hi-taxis-and-tyres}
\hlchaptermodality{406}

\begin{chapteropening}
\textbf{By the end of this chapter:} \emph{I can take a taxi, deal with a flat tyre and carry a licence and a helmet.}

Complete the last lesson of chapter 406: I can take a taxi, deal with a flat tyre and carry a licence and a helmet.
\end{chapteropening}

\section[\texorpdfstring{\hi{टैक्सी} (ṭaiksī)}{टैक्सी (ṭaiksī)}]{\hi{टैक्सी} (ṭaiksī) — a taxi}

\label{lesson:HI-C406-taiksi}



\begin{quote}
Before the new one: say the Hindi for to bounce, then the Hindi for to be cut.
\end{quote}

\subsection*{You'll want to know: \hi{टैक्सी}}
\textbf{\hi{टैक्सी}} — \emph{ṭaiksī} — \textquotedblleft{}a taxi\textquotedblright{}.

\begin{quote}
\hi{हमने} \hi{स्टेशन} \hi{से} \hi{टैक्सी} \hi{ली।}
\end{quote}

\begin{quote}
\emph{hamne sṭeśan se ṭaiksī lī.}
\end{quote}

\begin{quote}
We took a taxi from the station.
\end{quote}

Feminine.

\begin{cousinweb}
\textbf{\hi{टैक्सी}} (\emph{ṭaiksī}) is borrowed from English \emph{taxi}; it is feminine in Hindi, like \hi{गाड़ी}.
\end{cousinweb}

\begin{grammarlens}[title={What you've built}]
One more word for getting around.
\end{grammarlens}

\subsection*{Your turn}
\begin{itemize}
\raggedright
  \item \emph{Say it:} \emph{ṭaiksī}
  \item \emph{Say it:} \emph{ṭaiksī}, once more
  \item \emph{Say it:} \emph{hamne sṭeśan se ṭaiksī lī.} — We took a taxi from the station.
  \item \emph{Recall:} say the Hindi for to bounce, then the Hindi for to be cut, then say \emph{ṭaiksī} again
\end{itemize}

\begin{tcolorbox}[breakable,colback=teal!4,colframe=teal!35!black,title={Before you move on}]
What does \hi{टैक्सी} mean? (\textquotedblleft{}A taxi\textquotedblright{}.) Say it once more.
\end{tcolorbox}
\section[\texorpdfstring{\hi{टायर} (ṭāyar)}{टायर (ṭāyar)}]{\hi{टायर} (ṭāyar) — a tyre}

\label{lesson:HI-C406-tayar}



\begin{quote}
Before the new one: say the Hindi for to be cut, then the Hindi for a taxi.
\end{quote}

\subsection*{You'll want to know: \hi{टायर}}
\textbf{\hi{टायर}} — \emph{ṭāyar} — \textquotedblleft{}a tyre\textquotedblright{}.

\begin{quote}
\hi{गाड़ी} \hi{का} \hi{एक} \hi{टायर} \hi{पुराना} \hi{था।}
\end{quote}

\begin{quote}
\emph{gāṛī kā ek ṭāyar purānā thā.}
\end{quote}

\begin{quote}
One of the car's tyres was old.
\end{quote}

Masculine.

\begin{cousinweb}
\textbf{\hi{टायर}} (\emph{ṭāyar}) is borrowed from English \emph{tyre}, the rubber round a wheel.
\end{cousinweb}

\begin{grammarlens}[title={What you've built}]
One more word for getting around.
\end{grammarlens}

\subsection*{Your turn}
\begin{itemize}
\raggedright
  \item \emph{Say it:} \emph{ṭāyar}
  \item \emph{Say it:} \emph{ṭāyar}, once more
  \item \emph{Say it:} \emph{gāṛī kā ek ṭāyar purānā thā.} — One of the car's tyres was old.
  \item \emph{Recall:} say the Hindi for to be cut, then the Hindi for a taxi, then say \emph{ṭāyar} again
\end{itemize}

\begin{tcolorbox}[breakable,colback=teal!4,colframe=teal!35!black,title={Before you move on}]
What does \hi{टायर} mean? (\textquotedblleft{}A tyre\textquotedblright{}.) Say it once more.
\end{tcolorbox}
\section[\texorpdfstring{\hi{पंचर} (pañcar)}{पंचर (pañcar)}]{\hi{पंचर} (pañcar) — a puncture; punctured}

\label{lesson:HI-C406-pancar}



\begin{quote}
Before the new one: say the Hindi for a taxi, then the Hindi for a tyre.
\end{quote}

\subsection*{You'll want to know: \hi{पंचर}}
\textbf{\hi{पंचर}} — \emph{pañcar} — \textquotedblleft{}a puncture; punctured\textquotedblright{}.

\begin{quote}
\hi{रास्ते} \hi{में} \hi{टायर} \hi{पंचर} \hi{हो} \hi{गया।}
\end{quote}

\begin{quote}
\emph{rāste meṁ ṭāyar pañcar ho gayā.}
\end{quote}

\begin{quote}
The tyre got a puncture on the way.
\end{quote}

Masculine.

\begin{cousinweb}
\textbf{\hi{पंचर}} (\emph{pañcar}) is borrowed from English \emph{puncture}, shortened in the mouth; \hi{पंचर} \hi{होना}, to get a flat tyre.
\end{cousinweb}

\begin{grammarlens}[title={What you've built}]
One more word for getting around.
\end{grammarlens}

\subsection*{Your turn}
\begin{itemize}
\raggedright
  \item \emph{Say it:} \emph{pañcar}
  \item \emph{Say it:} \emph{pañcar}, once more
  \item \emph{Say it:} \emph{rāste meṁ ṭāyar pañcar ho gayā.} — The tyre got a puncture on the way.
  \item \emph{Recall:} say the Hindi for a taxi, then the Hindi for a tyre, then say \emph{pañcar} again
\end{itemize}

\begin{tcolorbox}[breakable,colback=teal!4,colframe=teal!35!black,title={Before you move on}]
What does \hi{पंचर} mean? (\textquotedblleft{}A puncture; punctured\textquotedblright{}.) Say it once more.
\end{tcolorbox}
\section[\texorpdfstring{\hi{हेलमेट} (helmeṭ)}{हेलमेट (helmeṭ)}]{\hi{हेलमेट} (helmeṭ) — a helmet}

\label{lesson:HI-C406-helmet}



\begin{quote}
Before the new one: say the Hindi for a tyre, then the Hindi for a puncture.
\end{quote}

\subsection*{You'll want to know: \hi{हेलमेट}}
\textbf{\hi{हेलमेट}} — \emph{helmeṭ} — \textquotedblleft{}a helmet\textquotedblright{}.

\begin{quote}
\hi{मोटरसाइकिल} \hi{पर} \hi{हमेशा} \hi{हेलमेट} \hi{पहनिए।}
\end{quote}

\begin{quote}
\emph{moṭarsāikil par hameśā helmeṭ pahaniye.}
\end{quote}

\begin{quote}
Always wear a helmet on a motorcycle.
\end{quote}

Masculine.

\begin{cousinweb}
\textbf{\hi{हेलमेट}} (\emph{helmeṭ}) is borrowed from English \emph{helmet}.
\end{cousinweb}

\begin{grammarlens}[title={What you've built}]
One more word for getting around.
\end{grammarlens}

\subsection*{Your turn}
\begin{itemize}
\raggedright
  \item \emph{Say it:} \emph{helmeṭ}
  \item \emph{Say it:} \emph{helmeṭ}, once more
  \item \emph{Say it:} \emph{moṭarsāikil par hameśā helmeṭ pahaniye.} — Always wear a helmet on a motorcycle.
  \item \emph{Recall:} say the Hindi for a tyre, then the Hindi for a puncture, then say \emph{helmeṭ} again
\end{itemize}

\begin{tcolorbox}[breakable,colback=teal!4,colframe=teal!35!black,title={Before you move on}]
What does \hi{हेलमेट} mean? (\textquotedblleft{}A helmet\textquotedblright{}.) Say it once more.
\end{tcolorbox}
\section[\texorpdfstring{\hi{लाइसेंस} (lāisens)}{लाइसेंस (lāisens)}]{\hi{लाइसेंस} (lāisens) — a licence}

\label{lesson:HI-C406-laisens}



\begin{quote}
Before the new one: say the Hindi for a puncture, then the Hindi for a helmet.
\end{quote}

\subsection*{You'll want to know: \hi{लाइसेंस}}
\textbf{\hi{लाइसेंस}} — \emph{lāisens} — \textquotedblleft{}a licence\textquotedblright{}.

\begin{quote}
\hi{ड्राइवर} \hi{के} \hi{पास} \hi{लाइसेंस} \hi{नहीं} \hi{था।}
\end{quote}

\begin{quote}
\emph{ḍrāivar ke pās lāisens nahīṃ thā.}
\end{quote}

\begin{quote}
The driver did not have a licence.
\end{quote}

Masculine.

\begin{cousinweb}
\textbf{\hi{लाइसेंस}} (\emph{lāisens}) is borrowed from English \emph{licence}, a driving licence above all.
\end{cousinweb}

\begin{grammarlens}[title={What you've built}]
One more word for getting around.
\end{grammarlens}

\subsection*{Your turn}
\begin{itemize}
\raggedright
  \item \emph{Say it:} \emph{lāisens}
  \item \emph{Say it:} \emph{lāisens}, once more
  \item \emph{Say it:} \emph{ḍrāivar ke pās lāisens nahīṃ thā.} — The driver did not have a licence.
  \item \emph{Recall:} say the Hindi for a puncture, then the Hindi for a helmet, then say \emph{lāisens} again
\end{itemize}

\begin{tcolorbox}[breakable,colback=teal!4,colframe=teal!35!black,title={Before you move on}]
What does \hi{लाइसेंस} mean? (\textquotedblleft{}A licence\textquotedblright{}.) Say it once more.
\end{tcolorbox}
